\begin{abstract}

  Large Language Models (LLMs) are increasingly deployed in safety-critical
  domains to generate structured outputs, but they struggle to satisfy complex,
  context-sensitive semantic invariants. While constrained decoding can enforce
  context-free grammars, it cannot guarantee richer typed or state-dependent
  constraints without incurring large inference overheads.
  We propose a dual-objective neuro-symbolic framework that integrates 
  runtime formal verification with structural fine-tuning. A deterministic
  \textit{Syntax Oracle}, an abstract verification machine, will maintain 
  specification state during decoding, and projects invalid proposals onto the
  admissable constraint spaec via dynamic logit masking. This mechanism provides
  formal structural guarantees, while acting only as a minimal-intervention
  safety controller.
  In parallel, we introduce a semantic internalization pipeline using
  parameter-efficient fine-tuning with Low-Rank Adaptation (LoRA), training the
  LLM to internalize structural invariants and thereby reduce reliance on
  deterministic correction.
  We evaluate the framework using an entropy-normalized Oracle Intervention
  Rate, alongside latency and throughput benchmarks. This approach aims to
  provide formally-verified, low-latency structured generation, while
  quantifying the trade-off between learned structural competence and runtime
  intervention.

\end{abstract}

\begin{IEEEkeywords}
  Neural networks, Language models, Formal verification, Runtime verification,
  Programming languages
\end{IEEEkeywords}


